\section{System Model}
The system is as follows.

Let $\mathcal{N} = \{1, 2, \dots, N\}$ denote the set of ground robots. Each robot $i \in \mathcal{N}$ generates a computationally intensive task modeled as a tuple $D_i = \{s_i, c_i, \tau_i\}$, where $s_i$ is the input data size in bits, $c_i$ is the required CPU cycles per bit, and $\tau_i$ is the maximum tolerable latency in seconds.

The system relies on the following decision variables:
\begin{itemize}
    \item $x_i \in \{0, 1\}$: The offloading decision indicator. $x_i = 1$ if task $i$ is processed locally at the Drone-AP; $x_i = 0$ if the task is relayed to the Cloud.
    \item $P_i \ge 0$: The transmit power of robot $i$ on the Ground-to-Air (G2A) link.
    \item $f_{D,i} \ge 0$: The Drone CPU frequency allocated to task $i$.
    \item $f_{C,i} \ge 0$: The Cloud CPU frequency allocated to task $i$.
\end{itemize}
\subsection{Ground-to-Air (G2A) Channel}
The G2A channel between the ground robots and the Drone-AP is modeled using a probabilistic Line-of-Sight (LoS) and Non-Line-of-Sight (NLoS) framework. The path loss for robot $i$, denoted $L^{\text{G2A}}_i$ in dB, is expressed as:
\begin{equation}
L^{\text{G2A}}_i = 20 \log_{10}\left(\frac{4\pi f_c d_i}{c}\right) + \eta
\end{equation}
where $d_i$ is the 3D distance between robot $i$ and the drone, $f_c$ is the carrier frequency, $c$ is the speed of light, and $\eta$ is the additional attenuation factor depending on whether the link is LoS ($\eta_{\text{LoS}}$) or NLoS ($\eta_{\text{NLoS}}$).\hl{Add Al-Hourani's channel model here!!}

The achievable uplink data rate is given by:
\begin{equation}
R^{\text{G2A}}_i = B_{\text{G2A}} \log_2 \left( 1 + \frac{P_i 10^{-L^{\text{G2A}}_i/10}}{N_0 B_{\text{G2A}}} \right)
\end{equation}
where $B_{\text{G2A}}$ is the channel bandwidth and $N_0$ is the noise power spectral density.\hl{take the channel bw value from simulator!}

\subsection{Air-to-Space (A2S) Channel}
The backhaul link from the Drone-AP to the LEO satellite relies on a strong LoS channel. The path loss $L^{\text{A2S}}$ incorporates Free Space Path Loss and atmospheric attenuation $\xi_{\text{atm}}$:
\begin{equation}
L^{\text{A2S}} = 20 \log_{10}\left(\frac{4\pi f_{\text{sat}} d^{t}_{\text{sat}}}{c}\right) + \xi_{\text{atm}},
\end{equation}
where $d^{t}_{\text{sat}}$ indicates the distance between the satellite and the UAV-BS at time $t$.\hl{consider using $\bar{d}$ to indicate average d while uav and sat are associated!}

The achievable A2S data rate for relayed tasks is:
\begin{equation}
R^{\text{A2S}}_i = B_{\text{A2S}} \log_2 \left( 1 + \frac{P_D |g|^2 10^{-L^{\text{A2S}}/10}}{N_0 B_{\text{A2S}}} \right)
\end{equation}
where $P_D$ is the transmission power of the Drone-AP and $g$ models the Rician fading component.\hl{find the bw value and exact channel parameters!}

\hl{We need to add space-air and air-to-ground downlink parameters. Only changing the tx power may be sufficient!}

\subsection{Latency Model}
The latency experienced by task $i$ depends on the offloading decision. If computed at the Drone-AP ($x_i = 1$):
\begin{equation}
T_{i, \text{Drone}} = \frac{s_i}{R^{\text{G2A}}_i} + \frac{s_i c_i}{f_{D,i}}
\end{equation}

If the task is relayed to the Cloud ($x_i = 0$):
\begin{equation}
T_{i, \text{Cloud}} = \frac{s_i}{R^{\text{G2A}}_i} + \frac{s_i}{R^{\text{A2S}}_i} + T_{\text{prop}} + \frac{s_i c_i}{f_{C,i}}
\end{equation}
where $T_{\text{prop}} \approx 2 \bar{d}_{\text{sat}} / c$ represents the satellite round-trip propagation delay by using the average of the distance between the UAV and the satellite while associated. The total latency for task $i$ is thus:
\begin{equation}
T_i = x_i T_{i, \text{Drone}} + (1 - x_i) T_{i, \text{Cloud}}
\end{equation}

\subsection{Energy Consumption Model}
The transmission energy consumed by robot $i$ is:
\begin{equation}
E^{\text{tx}}_i = P_i \frac{s_i}{R^{\text{G2A}}_i}
\end{equation}\hl{confirm that this needs to be fixed for 3D placement}
The energy expended by the Drone-AP encompasses local computation and relaying transmissions. Let $\kappa_D$ denote the \hl{effective switched capacitance} of the drone's computing architecture\hl{is this well known and widely used? how about GOP?}. The drone energy cost for task $i$ is:
\begin{equation}
E_{\text{Drone}, i} = x_i \kappa_D (f_{D,i})^2 s_i c_i + (1 - x_i) P_D \frac{s_i}{R^{\text{A2S}}_i}
\end{equation}
% BEGIN PANAMA ADDITIONS: communication and information state
\begin{panamaaddition}
\subsection{Timestamped Communication and Information Age}
Let $n\in\mathbb{N}_0$ index joint robot actions, with decision time
$t_n=n\Delta$ and step duration $\Delta>0$. Robot $i$ occupies grid position
$\mathbf{q}_i(n)$ and has a known goal $\mathbf{g}_i$. Its physical state is
distinct from the information stored at a remote processor or another robot.
Each message $m$ carries its source timestamp $g_m$, sequence or model version
$\nu_m$, and complete payload. Its receiver obtains that payload at
\begin{equation}
 b_m=g_m+\delta_m,\qquad \delta_m\geq0,
 \label{eq:panama-arrival}
\end{equation}
where $\delta_m$ is the end-to-end delay supplied by the communication
transport. A message with no completed delivery has $b_m=+\infty$. The
algorithm does not prescribe a distribution for $\delta_m$ or a handover
schedule. A subsequent LEO model determines these arrival times.

For a receiver stream $f$, let $\mathcal{M}_i^f(t_n)$ contain its messages
fully received by $t_n$. We consider task guidance $f=G$, peer intentions
$f=I_j$, and peer priorities $f=P_j$, for $j\in\mathcal{N}$. The accepted
message and corresponding cache are
\begin{align}
 m_i^{f,*}(n)&=\operatorname*{arg\,max}_{m\in\mathcal{M}_i^f(t_n)}
                         (g_m,\nu_m),\label{eq:panama-cache}\\
 \widehat{z}_i^f(n)&=z_{m_i^{f,*}(n)},
\end{align}
where the maximum uses lexicographic timestamp--version order. An initial
mission briefing initializes every cache at $t_0$. If no newer message
arrives, the stored value persists. An out-of-order older message cannot
overwrite a newer accepted cache.

A relayed or aggregated payload also retains the timestamp of each source
report, denoted $s_m^f$. Its information age is
\begin{equation}
 A_i^f(n)=t_n-s_{m_i^{f,*}(n)}^f.
 \label{eq:panama-age}
\end{equation}
For guidance, this is the timestamp of the robot state used to compute the
route. For intentions and priorities, it is the timestamp of the relevant
peer report. A newly received message can therefore contain old source
information. Aggregated peer age is the maximum of the included source ages.

Robots communicate through the UAV. For peer $j$ and receiver $i$, the
transport delay $\delta_{j\rightarrow i}^{f}$ covers the report and returned
update along the selected processing route. UAV processing permits local
exchange without the satellite backhaul, whereas cloud processing requires
backhaul forwarding. A cutoff prevents new deliveries on the affected route
and increases cache age. It does not remove the last received peer intention
or replace the robot's known goal. UAV positioning uses its own latest
received robot-position reports, $\widehat{\mathbf{q}}_i^U(t)$, rather than
instantaneous global positions unavailable to that receiver.

\subsection{Physical State and Available Decision Information}
Self-position, locally sensed obstacles, and visible robot positions remain
current at each action. Communication affects only the designated cached
guidance, intentions, or priorities. The simulator resolves physical
collisions, calculates task rewards, and checks mission completion using the
true state. Cached priorities modify observation construction, not the
physical transition rule. An age value is recorded for analysis; it is not
an additional input to the retained policy. Thus the delay extension preserves
the policy input dimensions while changing the timestamps of selected inputs.
\end{panamaaddition}
% END PANAMA ADDITIONS: communication and information state
